\slajdid{09_3-s066}
\begin{frame}[label=09_3-s066]
\gusttekst\setlength{\abovedisplayskip}{3pt}\setlength{\belowdisplayskip}{3pt}\setlength{\abovedisplayshortskip}{2pt}\setlength{\belowdisplayshortskip}{2pt}Situacija kada se prekidač koristi kao transmisioni gejt u digitalnoj logici je potpuno identična. Ako je kontrolni signal jednak 1 na N tranzistoru i 0 na P tranzistoru oba tranzistora „imaju uslove za provođenje“ u zavisnosti od ulaznog napona. Ulazni napon 1, vodi P tranzistor, ulazni napon 0 vodi N tranzistor. Ako je kontrolni signal jednak 0 na N tranzistoru i 1 na P tranzistoru, oba tranzistora su zakočena.
\begin{center}\begin{tikzpicture}[x=.85cm,y=.6cm,font=\sitneoznake,>=Latex]\begin{scope}[shift={(0,1.5)}]\draw(-.3,.4)--(.3,.4);\draw(-.4,.53)--(.4,.53);\draw(-.3,-.4)--(.3,-.4);\draw(-.4,-.53)--(.4,-.53);\draw(-.3,.4)--(-.3,-.4);\draw(.3,.4)--(.3,-.4);\draw(0,.65)circle(.12);\draw(0,.77)--(0,1.1);\draw(0,-.53)--(0,-1.1);\draw(-1,0)node[ocirc]{}node[left]{A}--(-.3,0);\draw(.3,0)--(1.5,0);\end{scope}\begin{scope}[shift={(0,-1.5)}]\draw(-.3,.4)--(.3,.4);\draw(-.4,.53)--(.4,.53);\draw(-.3,-.4)--(.3,-.4);\draw(-.4,-.53)--(.4,-.53);\draw(-.3,.4)--(-.3,-.4);\draw(.3,.4)--(.3,-.4);\draw(0,.65)circle(.12);\draw(0,.77)--(0,1.1);\draw(0,-.53)--(0,-1.1);\draw(-1,0)node[ocirc]{}node[left]{B}--(-.3,0);\draw(.3,0)--(1.5,0);\end{scope}\node[above]at(0,2.6){$\overline S$};\node[below]at(0,-2.6){$\overline S$};\draw(0,.4)--(0,-.4);\draw(-1,.1)node[ocirc]{}node[left]{$S$}--(0,.1);
\draw(1.5,1.5)--(1.5,-1.5);\draw(1.5,0)--(2.3,0);\draw(2.3,-1.5)--(2.3,1.5);
\draw(3.15,1.1)--(3.15,1.9);\draw(3,1.05)--(3,1.95);\draw(2.3,1.5)--(2.76,1.5);\draw(2.88,1.5)circle(.12);\draw(3.15,1.85)--(3.6,1.85)--(3.6,2.6);\draw(3.3,2.6)--(3.9,2.6);\node[above]at(3.6,2.6){$V_{DD}$};\draw(3.15,1.15)--(3.6,1.15)--(3.6,-1.15);
\draw(3.15,-1.9)--(3.15,-1.1);\draw(3,-1.95)--(3,-1.05);\draw(2.3,-1.5)--(3,-1.5);\draw(3.15,-1.15)--(3.6,-1.15);\draw(3.15,-1.85)--(3.6,-1.85)--(3.6,-2.6)node[ground]{};
\draw(3.6,0)--(4.3,0)node[ocirc]{}node[right]{$\overline Y=SA+\overline S B$};\end{tikzpicture}
\end{center}
\end{frame}
